\chapter{データ分析と可視化}
\label{chap7}

\begin{chapterbox}
\textbf{【この章で学ぶこと】}
\begin{itemize}
\item AIを活用したデータ分析ワークフロー（前処理→分析→可視化）を実践できる
\item 論文品質のグラフを効率的に作成できる
\item 統計的検定の基礎と適切な手法選択を理解する
\end{itemize}
\end{chapterbox}
\vspace{5mm}

\section{研究におけるデータ分析}
\index{でーたぶんせき@データ分析}

\subsection{なぜデータ分析が重要か}

大学院での研究において、データ分析\index{でーたぶんせき@データ分析}は以下の場面で欠かせない。

\begin{enumerate}
\item \textbf{実験結果の解釈}: 収集したデータから意味のあるパターンを見出す
\item \textbf{仮説の検証}: 統計的手法により、仮説が支持されるかを客観的に判断する
\item \textbf{論文の説得力}: 適切な分析と可視化は、論文の査読者を説得する論文の説得力を大きく左右する
\item \textbf{新たな発見}: 探索的データ分析（EDA）により、予想外の知見が得られることがある
\end{enumerate}

データ分析のスキルは、理系・文系を問わず、現代の大学院生に求められる基本能力である。心理学の実験データ、社会調査の回答データ、生物学の測定データ、工学のシミュレーション結果——あらゆる研究分野でデータ分析が行われている。

\subsection{AIがデータ分析をどう変えるか}

従来、データ分析を行うには、プログラミング言語（Python, R等）の習得に多くの時間を要した。AIの登場により、このハードルは大幅に下がっている。

\begin{itembox}[l]{AIが特に役立つ場面}
\begin{itemize}
\item \textbf{分析コードの生成}: 「このデータでt検定を行いたい」と伝えるだけで、コードが得られる
\item \textbf{手法の選択支援}: データの特性を伝えると、適切な統計手法を提案してくれる
\item \textbf{グラフの作成}: 「論文に載せる棒グラフを作って」という指示で、出版品質のグラフコードが得られる
\item \textbf{結果の解釈}: 統計結果の意味を平易な言葉で説明してもらえる
\end{itemize}
\end{itembox}

ただし、\textbf{AIはあくまで道具であり、分析の設計と結果の解釈は研究者の責任である}。本章では、AIの力を正しく活用しつつ、データ分析の基礎力を養うことを目指す。

\section{データ分析ワークフロー}
\index{わーくふろー@ワークフロー}

\subsection{Step 1: データの理解と前処理}

データ分析の第一歩は、手元のデータを正しく理解することである。

\begin{promptbox}
以下のCSVファイルを読み込み、データの概要を把握する\\
Pythonスクリプトを作成してください。

ファイル: experiment\_results.csv\\
使用ライブラリ: pandas

以下の情報を出力してください:
\begin{enumerate}
\item データの形状（行数×列数）
\item 各列のデータ型
\item 最初の5行の表示
\item 基本統計量（数値列: 平均・標準偏差・最小・最大、\\
カテゴリ列: ユニーク数・最頻値）
\item 欠損値の数と割合（列ごと）
\item 各数値列のヒストグラム
\end{enumerate}
\end{promptbox}

AIが生成するコードの例:

\begin{lstlisting}[language=Python]
import pandas as pd
import matplotlib.pyplot as plt
import numpy as np


def explore_data(file_path: str) -> pd.DataFrame:
    """CSVファイルを読み込み、データの概要を表示する"""

    # データ読み込み
    df = pd.read_csv(file_path)

    # 1. データの形状
    print("=" * 50)
    print(f"データの形状: {df.shape[0]}行 × {df.shape[1]}列")
    print("=" * 50)

    # 2. 各列のデータ型
    print("\n【データ型】")
    print(df.dtypes.to_string())

    # 3. 最初の5行
    print("\n【最初の5行】")
    print(df.head().to_string())

    # 4. 基本統計量
    print("\n【基本統計量（数値列）】")
    print(df.describe().to_string())

    # 5. 欠損値
    print("\n【欠損値】")
    missing = df.isnull().sum()
    missing_pct = (missing / len(df) * 100).round(2)
    missing_df = pd.DataFrame({
        "欠損数": missing,
        "欠損率(%)": missing_pct
    })
    print(missing_df[missing_df["欠損数"] > 0].to_string())
    if missing.sum() == 0:
        print("  欠損値なし")

    # 6. ヒストグラム
    num_cols = df.select_dtypes(include=[np.number]).columns
    if len(num_cols) > 0:
        fig, axes = plt.subplots(
            1, len(num_cols),
            figsize=(5 * len(num_cols), 4))
        if len(num_cols) == 1:
            axes = [axes]
        for ax, col in zip(axes, num_cols):
            df[col].hist(ax=ax, bins=20, edgecolor="black")
            ax.set_title(col)
            ax.set_xlabel("値")
            ax.set_ylabel("頻度")
        plt.tight_layout()
        plt.savefig("data_overview_histograms.png", dpi=150)
        plt.show()

    return df
\end{lstlisting}

\begin{itembox}[l]{欠損値処理の主な方法}
\begin{table}[H]
\centering
\footnotesize
\begin{tabular}{lll}
\toprule
方法 & 適用場面 & Pythonコード \\
\midrule
行の削除 & 欠損が少ない場合 & \texttt{df.dropna()} \\
平均値補完 & 正規分布に近いデータ & \texttt{df.fillna(df.mean())} \\
中央値補完 & 外れ値がある場合 & \texttt{df.fillna(df.median())} \\
前方補完 & 時系列データ & \texttt{df.ffill()} \\
多重代入法 & 欠損が多い場合 & \texttt{IterativeImputer} \\
\bottomrule
\end{tabular}
\end{table}
\end{itembox}

\begin{promptbox}
以下のDataFrameの数値列について、外れ値を検出する\\
コードを作成してください。\\
外れ値の定義: IQR法（Q1 - 1.5*IQR 未満、\\
または Q3 + 1.5*IQR 超過）\\
外れ値がある行を表示し、箱ひげ図で可視化してください。
\end{promptbox}

\subsection{Step 2: 探索的データ分析（EDA）}
\index{たんさくてきでーたぶんせき@探索的データ分析}

探索的データ分析（Exploratory Data Analysis, EDA）\index{EDA@EDA}は、データの構造やパターンを視覚的に把握するプロセスである。

\begin{promptbox}
以下のDataFrameについて、探索的データ分析（EDA）を行う\\
Pythonコードを作成してください。

分析内容:
\begin{enumerate}
\item 数値列間の相関行列の計算とヒートマップの作成
\item 各数値列の分布（ヒストグラム + 正規分布フィット）
\item カテゴリ列ごとのグループ別統計量
\item 主要な変数間の散布図マトリクス
\end{enumerate}

使用ライブラリ: pandas, matplotlib, seaborn\\
グラフは日本語フォントに対応させてください。
\end{promptbox}

\subsection{Step 3: 統計的検定}
\index{とうけいてきけんてい@統計的検定}

研究の仮説を検証するためには、適切な統計的検定を行う必要がある。以下に、よく使われる検定手法とその選び方を示す。

\begin{itembox}[l]{統計的検定の選び方フローチャート}
\begin{verbatim}
データは正規分布に従うか？
├── はい → 2群比較か？
│   ├── はい → 対応のあるデータか？
│   │   ├── はい → 対応のあるt検定
│   │   └── いいえ → 独立2標本t検定
│   └── いいえ → 3群以上の比較
│       ├── 対応あり → 反復測定分散分析
│       └── 対応なし → 一元配置分散分析
└── いいえ → 2群比較か？
    ├── はい → 対応のあるデータか？
    │   ├── はい → Wilcoxon符号順位検定
    │   └── いいえ → Mann-Whitney U検定
    └── いいえ → Kruskal-Wallis検定
カテゴリデータ同士 → カイ二乗検定
\end{verbatim}
\end{itembox}

\begin{promptbox}
以下のデータについて、2群間の平均値の差を検定する\\
Pythonコードを作成してください。

データ構造:
\begin{itemize}
\item DataFrame に ``group'' 列（``control'' と ``treatment''）と ``score'' 列がある
\item 各群の被験者数は異なる可能性がある
\end{itemize}

処理手順:
\begin{enumerate}
\item 各群のデータの正規性を Shapiro-Wilk 検定で確認
\item 等分散性を Levene 検定で確認
\item 結果に応じて適切な検定を自動選択
\item 効果量（Cohen's d または r）を計算
\item 結果を分かりやすくレポート出力
\end{enumerate}
\end{promptbox}

\begin{itembox}[l]{p値の正しい解釈}
\index{pち@$p$値}
\begin{table}[H]
\centering
\begin{tabular}{p{4.5cm}p{5cm}}
\toprule
誤った解釈 & 正しい解釈 \\
\midrule
$p = 0.03$は「差がある確率が97\%」 & 帰無仮説の下で、観測データ以上に極端な結果が得られる確率が3\% \\
$p < 0.05$は「重要な差がある」 & 統計的に有意だが、実用的な重要性は効果量で判断 \\
$p = 0.06$は「差がない」 & 有意水準を超えただけで、差がないとは限らない \\
$p$値が小さいほど効果が大きい & $p$値はサンプルサイズの影響を受ける。効果量を見るべき \\
\bottomrule
\end{tabular}
\end{table}
\end{itembox}

\section{論文品質のグラフ作成}
\index{ぐらふさくせい@グラフ作成}

\subsection{matplotlib / seaborn の基本}

Pythonでグラフを作成するための2大ライブラリがmatplotlib\index{matplotlib@matplotlib}とseaborn\index{seaborn@seaborn}である。

\begin{itemize}
\item \textbf{matplotlib}: 低レベルの描画ライブラリ。細かいカスタマイズが可能
\item \textbf{seaborn}: matplotlibをベースにした高レベルの可視化ライブラリ。統計的な可視化に特化し、美しいデフォルト設定を持つ
\end{itemize}

\subsection{効果的なグラフの選び方}

\begin{table}[H]
\centering
\small
\begin{tabular}{lll}
\toprule
データの性質 & 推奨グラフ & 用途 \\
\midrule
カテゴリ別の量的比較 & 棒グラフ（エラーバー付き） & 群間比較の結果表示 \\
2変数の関係 & 散布図 & 相関関係の可視化 \\
分布の比較 & 箱ひげ図 / バイオリンプロット & 群間の分布の違い \\
変数間の関連 & ヒートマップ & 相関行列の可視化 \\
時系列データ & 折れ線グラフ & 時間変化の表示 \\
割合・構成比 & 積み上げ棒グラフ & カテゴリの内訳表示 \\
\bottomrule
\end{tabular}
\end{table}

\subsection{論文に必要なグラフの要素}

学術論文や学会発表のグラフには、以下の要素が不可欠である。

\begin{itembox}[l]{必須要素}
\begin{itemize}
\item 明確な軸ラベル（単位を含む）
\item 適切なタイトル（Figure 1: ... の形式）
\item 凡例（複数系列の場合）
\item エラーバー（平均値を示す場合、標準誤差または95\%信頼区間）
\item 適切なフォントサイズ（本文と同程度、最低8pt）
\item 高解像度（300dpi以上）
\end{itemize}
\end{itembox}

\subsection{白黒印刷対応のグラフ設計}

多くの学術論文は白黒で印刷されることがあるため、色だけに頼らないデザインが重要である。

\begin{enumerate}
\item \textbf{ハッチパターン}: 棒グラフにストライプや格子模様を加える
\item \textbf{マーカーの形状}: 散布図で○、△、□など異なるマーカーを使う
\item \textbf{線のスタイル}: 折れ線グラフで実線、破線、点線を使い分ける
\item \textbf{コントラスト}: 色を使う場合も、白黒にしたときに区別できるコントラストを確保する
\end{enumerate}

\subsection{AIでグラフ作成コードを生成するプロンプト例}

\begin{promptbox}
以下の仕様で論文品質の棒グラフを作成する\\
Pythonコードを書いてください。

データ:
\begin{itemize}
\item 3群（Control, Treatment A, Treatment B）
\item 各群の平均値と標準誤差がある
\end{itemize}

グラフの要件:
\begin{itemize}
\item matplotlib + seaborn を使用
\item 白背景、枠線あり
\item エラーバー（標準誤差）付き
\item 各棒に異なるハッチパターン（白黒印刷対応）
\item 有意差を示すブラケットと * マーク
\item フォントサイズ: タイトル14pt、軸ラベル12pt、目盛り10pt
\item 解像度 300dpi で PNG 保存
\item 図のサイズ: 幅8cm × 高さ6cm（単カラム論文用）
\end{itemize}
\end{promptbox}

\section{インタラクティブ可視化}
\index{いんたらくてぃぶかしか@インタラクティブ可視化}

\subsection{Plotlyの紹介}

Plotly\index{Plotly@Plotly}は、インタラクティブなグラフを作成するライブラリである。マウスオーバーで値を表示したり、ズーム・パンができるため、データの探索段階で特に有用である。

\begin{lstlisting}[language=Python]
import plotly.express as px

# インタラクティブ散布図
fig = px.scatter(
    df, x="variable_a", y="variable_b",
    color="group",
    hover_data=["subject_id", "age", "score"],
    title="変数Aと変数Bの関係",
    labels={
        "variable_a": "変数A",
        "variable_b": "変数B"
    }
)
fig.update_layout(
    font=dict(size=14),
    hoverlabel=dict(font_size=12)
)
fig.write_html("interactive_scatter.html")
fig.show()
\end{lstlisting}

\subsection{データ探索での活用}

インタラクティブ可視化は、以下の場面で特に有用である。

\begin{itemize}
\item \textbf{外れ値の特定}: マウスオーバーで外れ値の詳細情報を確認できる
\item \textbf{パターンの発見}: ズーム機能で特定の領域を拡大して観察できる
\item \textbf{プレゼンテーション}: ゼミ発表でリアルタイムにデータを探索できる
\item \textbf{共同研究者との共有}: HTMLファイルを送るだけで、相手もインタラクティブに操作できる
\end{itemize}

ただし、\textbf{論文に掲載するグラフは静的な画像（matplotlib等）で作成する}こと。インタラクティブ可視化はあくまで探索・確認用である。

\section{AIを使ったデータ分析の落とし穴}

\subsection{統計的な誤解}
\index{pはっきんぐ@$p$ハッキング}

AIがデータ分析を手軽にする一方で、統計的な誤解に基づく分析が行われるリスクも増大している。

\begin{itembox}[l]{p-hacking（p値ハッキング）}
多数の検定を繰り返し、有意な結果が出たものだけを報告する行為。AIを使うと、様々な分析を手軽に試せるため、このリスクが高まる。

有意水準5\%で20回検定すれば、本当は差がなくても平均1回は「有意」になる（偽陽性）。
\end{itembox}

\begin{itembox}[l]{多重比較問題}
\index{たじゅうひかく@多重比較}
複数の群を同時に比較する場合、検定の回数が増えるほど偽陽性の確率が高くなる。3群を比較する場合はBonferroni補正（補正後の有意水準: $0.05 / 3 = 0.0167$）やTukey HSDなどの多重比較法を用いる。
\end{itembox}

\subsection{過学習への注意}
\index{かがくしゅう@過学習}

AIが提案する分析手法が必要以上に複雑になることがある。特に機械学習を使う場合、過学習（overfitting）のリスクに注意が必要である。

対策:
\begin{itemize}
\item 交差検証（Cross-validation）を必ず行う
\item 単純なモデルから始め、必要な場合にのみ複雑化する
\item サンプルサイズに見合った分析手法を選ぶ
\end{itemize}

\subsection{AIが出した統計結果を鵜呑みにしない}

AIが生成した分析コードは、一見正しく見えても論理的な誤りを含むことがある。

\begin{itembox}[l]{自分で確認すべきポイント}
\begin{itemize}
\item 検定の前提条件は満たされているか（正規性、等分散性）
\item 実験計画（対応あり/なし）に合った検定を使っているか
\item 有意水準は事前に設定されているか
\item 効果量も報告しているか
\item 多重比較を行う場合、補正は適切か
\item サンプルサイズは十分か
\item 結果の解釈は論理的か
\end{itemize}
\end{itembox}

\section*{演習問題}

\begin{enumerate}
\item[\textbf{演習7-1}] \textbf{データの前処理と概要把握}\\
サンプルデータを生成し、AIを使ってデータの概要把握と前処理を行いなさい。データの概要把握（形状、型、分布、欠損値）、外れ値の検出と対処、欠損値の処理、処理前後の比較を実施し、レポートにまとめなさい。

\item[\textbf{演習7-2}] \textbf{探索的データ分析}\\
演習7-1のデータについて、グループごとの記述統計量、相関分析、各変数の分布の可視化を実施し、発見した特徴やパターンを3つ以上報告しなさい。

\item[\textbf{演習7-3}] \textbf{適切な統計検定の実施}\\
「treatment群はcontrol群と比べて、pre\_scoreからpost\_scoreへの変化量が大きい」という仮説を検定しなさい。正規性の検定、適切な統計検定の選択と実施、効果量の計算を行い、結果を論文の「結果」セクションの書き方で報告しなさい。

\item[\textbf{演習7-4}] \textbf{論文品質のグラフ作成}\\
演習7-3の結果を以下のグラフで可視化しなさい。全てのグラフは論文品質（300dpi、適切なフォントサイズ、軸ラベル）で作成すること。(1) 群ごとのpre/postスコアの棒グラフ（エラーバー付き）、(2) 群ごとの変化量の箱ひげ図（個別データ点付き）、(3) pre\_scoreとpost\_scoreの散布図（群ごとに色分け、回帰直線付き）。

\item[\textbf{演習7-5}] \textbf{総合分析レポート}\\
自分の研究テーマについて、データ分析の全工程を実施し、レポートにまとめなさい。サンプルデータの生成、前処理とEDA、適切な統計分析、論文品質のグラフ3枚以上、結果のまとめ（発見した知見を3つ以上）を含めること。AIとの対話履歴の要約も合わせて提出すること。
\end{enumerate}
